\section{Tutorial 4: Regression vs Classification}

\subsection{Objectives}

This tutorial develops a clear understanding of the differences and connections between regression and classification:

\begin{itemize}
    \item Compare linear regression and logistic regression
    \item Understand loss functions and probabilistic interpretations
    \item Analyze decision boundaries
    \item Explore when each formulation is appropriate
\end{itemize}

\medskip

\noindent
\textbf{Focus.}  
Unify regression and classification under the framework of function approximation and probabilistic modeling.

\subsection{Exercise 1: Linear Regression for Classification}

Suppose we use linear regression for binary classification:
\[
f(x) = w^\top x, \quad \hat{y} = \text{sign}(f(x)).
\]

\medskip

\noindent
\textbf{Tasks.}
\begin{enumerate}
    \item What loss is being minimized?
    \item What are the potential issues with this approach?
\end{enumerate}

\medskip

\noindent
\textbf{Solution.}

\begin{enumerate}
\item The squared loss:
\[
\sum_i (w^\top x_i - y_i)^2.
\]

\item Issues:
\begin{itemize}
    \item Predictions are not probabilities
    \item Outputs can lie outside $[0,1]$
    \item Sensitive to outliers
    \item Does not align well with classification objectives
\end{itemize}
\end{enumerate}

\medskip

\noindent
\textbf{Key takeaway.}  
Regression loss is not well-suited for classification.

\subsection{Exercise 2: Logistic Regression Interpretation}

Recall:
\[
p(y=1 \mid x) = \sigma(w^\top x).
\]

\medskip

\noindent
\textbf{Tasks.}
\begin{enumerate}
    \item Why is the sigmoid function appropriate here?
    \item What does $w^\top x$ represent?
\end{enumerate}

\medskip

\noindent
\textbf{Solution.}

\begin{enumerate}
\item Sigmoid maps $\mathbb{R} \to (0,1)$, making it suitable for probabilities.

\item $w^\top x$ represents the log-odds:
\[
\log \frac{p(y=1 \mid x)}{p(y=0 \mid x)}.
\]
\end{enumerate}

\medskip

\noindent
\textbf{Key takeaway.}  
Logistic regression provides a principled probabilistic interpretation.

\subsection{Exercise 3: Loss Function Comparison}

Compare:

\begin{itemize}
    \item Squared loss (regression)
    \item Cross-entropy loss (classification)
\end{itemize}

\medskip

\noindent
\textbf{Tasks.}
\begin{enumerate}
    \item How do they penalize errors differently?
    \item Which is more appropriate for classification and why?
\end{enumerate}

\medskip

\noindent
\textbf{Solution.}

\begin{enumerate}
\item 
\begin{itemize}
    \item Squared loss penalizes distance symmetrically
    \item Cross-entropy penalizes confident wrong predictions heavily
\end{itemize}

\item Cross-entropy is more appropriate because:
\begin{itemize}
    \item It arises from likelihood
    \item It reflects uncertainty
    \item It aligns with classification objectives
\end{itemize}
\end{enumerate}

\medskip

\noindent
\textbf{Key takeaway.}  
Choice of loss reflects modeling assumptions.

\subsection{Exercise 4: Decision Boundaries}

Consider:
\[
f(x) = w^\top x.
\]

\medskip

\noindent
\textbf{Tasks.}
\begin{enumerate}
    \item Describe the decision boundary for regression (if thresholded).
    \item Describe the decision boundary for logistic regression.
\end{enumerate}

\medskip

\noindent
\textbf{Solution.}

\begin{enumerate}
\item Regression with thresholding:
\begin{itemize}
    \item Boundary: $w^\top x = c$ (depends on threshold)
\end{itemize}

\item Logistic regression:
\begin{itemize}
    \item Boundary: $w^\top x = 0$
\end{itemize}
\end{enumerate}

\medskip

\noindent
\textbf{Insight.}  
Both produce linear boundaries, but logistic regression has a probabilistic interpretation.

\subsection{Exercise 5: Effect of Feature Mapping}

Suppose we use features $\phi(x)$ instead of $x$.

\medskip

\noindent
\textbf{Tasks.}
\begin{enumerate}
    \item What happens to the decision boundary?
    \item How does this affect expressivity?
\end{enumerate}

\medskip

\noindent
\textbf{Solution.}

\begin{enumerate}
\item Decision boundary becomes:
\[
w^\top \phi(x) = 0,
\]
which can be nonlinear in the original input space.

\item Expressivity increases:
\begin{itemize}
    \item More complex patterns can be captured
\end{itemize}
\end{enumerate}

\medskip

\noindent
\textbf{Key takeaway.}  
Nonlinearity arises through representation, not necessarily through the model.

\subsection{Exercise 6 (Discussion): When to Use Regression vs Classification}

\textbf{Question.}  
When should we use regression, and when should we use classification?

\medskip

\noindent
\textbf{Discussion points.}

\begin{itemize}
    \item Regression:
    \begin{itemize}
        \item Continuous outputs
        \item Predicting quantities
    \end{itemize}

    \item Classification:
    \begin{itemize}
        \item Discrete labels
        \item Decision-making tasks
    \end{itemize}
\end{itemize}

\medskip

\noindent
\textbf{Subtle point.}
\begin{itemize}
    \item Classification models often estimate probabilities
    \item Final decisions arise from thresholds
\end{itemize}

\medskip

\noindent
\textbf{Insight.}  
Classification can be viewed as structured regression on probabilities.

\subsection{Summary}

\begin{itemize}
    \item Regression predicts continuous values; classification predicts labels
    \item Logistic regression provides probabilistic classification
    \item Cross-entropy aligns with likelihood-based modeling
    \item Both models can produce linear decision boundaries
    \item Nonlinearity arises through feature transformations
\end{itemize}

\medskip

\noindent
\textbf{Preparation for next lecture.}  
We now generalize linear models using feature maps and kernel methods to capture nonlinear relationships.